\section{Optimization-Based Control}

A central idea behind model predictive control is that the control action is not generated by a fixed pre-computed feedback law alone, but by repeatedly solving an optimization problem online. This point of view is best understood by comparing it with the classical control loop. In a standard feedback architecture, the controller receives a reference and the measured output of the plant, and then applies an input according to a control law that has been designed in advance. In optimization-based control, this fixed controller block is replaced, at least conceptually, by an optimizer.\\

The optimizer uses a model of the plant to predict how the system will evolve under candidate future inputs. Based on these predictions, it computes an input sequence that best satisfies the desired control objectives while respecting the relevant constraints. The first important shift is therefore methodological: instead of asking directly for a controller formula, we formulate a decision problem whose solution determines the input to be applied.\\
This approach is particularly natural when the control objective is tightly coupled with physical limitations and safety requirements. Consider, for instance, the problem of driving a race car around a track. One may wish to minimize the lap time, but this objective cannot be pursued arbitrarily. The car must stay on the road, avoid collisions, respect limitations on tire friction, and remain within admissible acceleration and steering capabilities. The control task is therefore inherently constrained. A meaningful strategy is to look ahead, predict the upcoming road geometry and the effect of future decisions, and then plan a feasible trajectory that achieves the best possible performance. This is precisely the kind of reasoning that optimization-based control formalizes.\\

The same idea can be described algorithmically as follows. At a given time instant, one solves an optimization problem to compute a planned trajectory and a corresponding sequence of future control actions. From this sequence, only the first control move is actually implemented. At the next sampling instant, the state is measured again, the prediction is updated, and a new optimization problem is solved. In this way, planning and feedback are combined in a single framework.

\begin{theorembox}[Optimization-based control]
The key idea is to compute the control input by solving, at each sampling instant, a finite-dimensional optimization problem built from three ingredients: a performance objective, a predictive model of the plant, and a set of constraints.
\end{theorembox}

\section{Concept of Model Predictive Control}

Model predictive control, usually abbreviated as MPC, is the most prominent realization of optimization-based control. The name itself summarizes its defining features. It is \emph{model-based} because future behavior is predicted through an internal model of the plant. It is \emph{predictive} because the controller explicitly reasons over a future horizon. Finally, it is a \emph{control} method because the resulting optimization problem is solved repeatedly in closed loop in order to decide the plant input.\\

A useful way to understand MPC is to compare it with more classical control design philosophies. In classical control, one typically designs a controller $C$ for a plant $P$ so as to achieve good disturbance rejection, limited sensitivity to measurement noise, and robustness to modeling uncertainty. The dominant language is often that of transfer functions, loop shaping, and frequency-domain specifications. By contrast, MPC is naturally expressed in the \textbf{time domain}. Its strength lies in the \textbf{explicit treatment of constraints}, both on the manipulated inputs and on the system variables. This makes it especially appealing in applications where safety, actuator saturation, and operating envelopes are of central importance.\\
Indeed, all physical systems are constrained. Inputs cannot be arbitrarily large because actuators have limits. Outputs and states are often restricted by performance or safety considerations, such as maximum temperature, pressure, velocity, or overshoot. In many practical settings, the economically or operationally optimal regime lies close to such limits. Traditional control methods often handle constraints only indirectly, for example by tuning the set-point conservatively or by adding ad hoc supervisory logic. MPC instead incorporates the constraints directly into the control design. This allows operation closer to the admissible boundaries, while still maintaining systematic decision-making.\\

Mathematically, MPC can be formulated over a finite prediction horizon of length $N$. At a generic time instant $k$, one considers the current state $x(k)$ and searches for a sequence of future control inputs
\[
U_k = \{u_k,u_{k+1},\dots,u_{k+N-1}\}
\]
that minimizes a cumulative performance criterion. A standard formulation is
\[
U_k^\star(x(k)) := \argmin_{U_k}\; \sum_{i=0}^{N-1} \stagecost(x_{k+i},u_{k+i})
\]
subject to the system dynamics
\[
x_{k+i+1}=Ax_{k+i}+Bu_{k+i},
\]
the initial condition
\[
x_k = x(k),
\]
and the constraints
\[
x_{k+i}\in\cX,\qquad u_{k+i}\in\cU,
\]
for all relevant indices $i$.\\
This compact expression already contains the essential ingredients of MPC. The objective function defines what the controller tries to optimize. The model provides predictions of future behavior. The constraint sets $\cX$ and $\cU$ encode admissible operating conditions. The decision variables are the future inputs across the prediction horizon.

The online operation of MPC follows a receding-horizon principle. At each sampling time, the current state is measured or estimated, the finite-horizon optimization problem is solved, and an optimal input sequence
\[
U_k^\star = \{u_k^\star,u_{k+1}^\star,\dots,u_{k+N-1}^\star\}
\]
is obtained. However, only the first element $u_k^\star$ is applied to the plant. After one sampling interval, the state is updated and the entire optimization is solved again from the new initial condition. The horizon therefore moves forward in time, which is why one also speaks of a \emph{receding horizon} strategy.\\
This repeated re-optimization is what \textbf{introduces feedback}. If the plant is disturbed, if the model is imperfect, or if the measured state differs from the previously predicted one, the next optimization step automatically adjusts the future plan. The controller does not blindly execute an open-loop plan; instead, it continuously replans using fresh information.

\begin{theorembox}[Receding horizon principle]
MPC computes an \emph{optimal sequence} of future control moves, but implements only the \emph{first} one. Feedback arises because the optimization problem is solved again at the next sampling instant using the newly measured or estimated state.
\end{theorembox}

\subsection{Why constraints matter}

The explicit inclusion of constraints is one of the main reasons for the success of MPC. Suppose one wants to regulate a process output to a desired set-point while an upper safety limit must never be violated. A classical controller may achieve fast convergence but produce overshoot that crosses the admissible bound. To avoid this, one may reduce aggressiveness or move the set-point away from the constraint, which typically leads to suboptimal operation. MPC offers a more systematic alternative: the constraint is included directly in the optimization problem, and the computed control action is chosen with full awareness of that limit.\\
This ability to reason ahead is crucial. The optimizer can anticipate that a constraint will become active in the near future and can begin adjusting the input before the violation occurs. In this sense, MPC is predictive not only in the use of a model, but also in the way it handles constraints proactively rather than reactively.

\section{Applications of MPC}

One of the remarkable features of MPC is the \textbf{breadth of its application domain}. The same conceptual framework appears in problems evolving at extremely different time scales.\\

At very fast time scales, optimization-based control ideas arise in computer and electronic systems, where decisions may need to be made on the \textbf{nanosecond or microsecond} level. Slightly slower, but still extremely fast, are power system applications, where control actions and operational decisions can be taken in the microsecond range depending on the problem formulation and hardware layer under consideration.\\
At the \textbf{millisecond} scale, traction control and automotive systems provide a classical application area. Here, one must coordinate performance objectives with safety and actuator limits in real time. On the scale of \textbf{seconds}, building control becomes relevant: thermal dynamics are slow enough to allow relatively long prediction horizons, yet constraints and energy-efficiency objectives play a major role. At the scale of \textbf{minutes and hours}, refinery operation and process control have historically been among the main industrial drivers for MPC. At even slower scales, \textbf{days or weeks}, one finds planning and scheduling problems such as train scheduling, nurse rostering, and production planning, where the same optimization-and-prediction philosophy remains useful.

These examples show that MPC should not be viewed only as a single algorithm for fast dynamical systems. Rather, it is a general methodology for repeated decision-making under dynamical models and constraints. The exact mathematical models, optimization problems, and computational requirements vary greatly across applications, but the underlying structure remains the same.

\section{A Brief History of MPC}

The historical roots of MPC go back to the early idea of using optimization methods directly inside control algorithms. An early milestone is the work of A.~I.~Propoi in 1963, where linear programming ideas were proposed for sampled-data automatic systems. This already contains the seed of a now-familiar idea: control can be posed as the repeated solution of an optimization problem.\\
The modern industrial development of MPC is commonly traced to the 1970s. A key step was the work of Richalet and co-authors in 1978 on what became known as IDCOM, short for \emph{Identification and Command}. Their approach used an impulse response model of the plant, considered a finite prediction horizon, and formulated a quadratic performance objective. Future plant behavior was described relative to a desired reference trajectory, and practical input/output constraints were included in an ad hoc manner. Because the resulting controller was not expressed as a conventional transfer function but rather as the outcome of an iterative computational procedure, the approach was described as heuristic predictive control. At this stage, the framework was mainly aimed at stable plants and was strongly motivated by industrial needs.\\

A closely related and highly influential development came from Cutler and collaborators. In the late 1970s, Cutler proposed ideas that were later implemented at Shell under the name \emph{Dynamic Matrix Control} (DMC). The 1979 work by Cutler and Ramaker is particularly important in the industrial history of MPC. DMC used a step response model of the plant, a finite-horizon quadratic objective, and a prediction mechanism that aimed at tracking the set-point as closely as possible. The optimal manipulated variables were computed through a least-squares problem, and additional equations were introduced online to account for constraints. The name ``dynamic matrix'' refers precisely to the structured matrix that arises in this least-squares formulation.\\
A further important step occurred in 1983, when Cutler, Morshedi, and Haydel emphasized the formulation of the controller as a standard quadratic program. This was conceptually decisive, because it provided a systematic mathematical framework for handling constraints rather than relying only on heuristic modifications. From that point on, MPC gradually evolved from an industrial engineering practice into a discipline with a firm optimization-theoretic foundation.\\

In the mid 1990s, much theoretical work was devoted to understanding under which conditions one can guarantee recursive feasibility and closed-loop stability. This period established many of the fundamental results that are now part of the standard theory of nominal MPC. During the 2000s, research expanded toward tractable robust MPC methods, nonlinear MPC, hybrid systems, and fast MPC algorithms tailored to increasingly demanding real-time requirements. In the 2010s, important developments included stochastic MPC, distributed and large-scale MPC, and economic MPC, where the objective is not merely set-point tracking but direct optimization of process performance. In the 2020s, a major direction has been the interaction between MPC and learning, leading to a growing body of work on learning-based MPC.

\section{Summary and Main Aspects}

Model predictive control can therefore be summarized as a control methodology in which, at every sampling instant, one measures or estimates the current state, predicts future behavior with an internal model, solves a constrained finite-horizon optimization problem, and applies only the first input of the optimal sequence. The procedure is then repeated at the next sampling instant.\\

Its main appeal lies in the fact that performance objectives and operating constraints are treated within a single optimization-based framework. This gives the designer a systematic way to handle actuator limits, safety requirements, and multi-variable interactions while maintaining a feedback structure through repeated re-planning.\\
At the same time, MPC also comes with \textbf{important challenges}. First, the optimization problem must be solved reliably and fast enough for real-time use. Second, the finite-horizon problem may become infeasible if the constraints are too restrictive or if the system is driven into an unfavorable region. Third, closed-loop stability is not automatic and must typically be enforced or analyzed through suitable design choices. Finally, the closed-loop behavior may be sensitive to model mismatch and disturbances, which motivates the study of robustness and the many extensions of the nominal MPC framework.
